% Options for packages loaded elsewhere
\PassOptionsToPackage{unicode}{hyperref}
\PassOptionsToPackage{hyphens}{url}
%
\documentclass[
]{article}
\usepackage{lmodern}
\usepackage{amssymb,amsmath}
\usepackage{ifxetex,ifluatex}
\ifnum 0\ifxetex 1\fi\ifluatex 1\fi=0 % if pdftex
  \usepackage[T1]{fontenc}
  \usepackage[utf8]{inputenc}
  \usepackage{textcomp} % provide euro and other symbols
\else % if luatex or xetex
  \usepackage{unicode-math}
  \defaultfontfeatures{Scale=MatchLowercase}
  \defaultfontfeatures[\rmfamily]{Ligatures=TeX,Scale=1}
\fi
% Use upquote if available, for straight quotes in verbatim environments
\IfFileExists{upquote.sty}{\usepackage{upquote}}{}
\IfFileExists{microtype.sty}{% use microtype if available
  \usepackage[]{microtype}
  \UseMicrotypeSet[protrusion]{basicmath} % disable protrusion for tt fonts
}{}
\makeatletter
\@ifundefined{KOMAClassName}{% if non-KOMA class
  \IfFileExists{parskip.sty}{%
    \usepackage{parskip}
  }{% else
    \setlength{\parindent}{0pt}
    \setlength{\parskip}{6pt plus 2pt minus 1pt}}
}{% if KOMA class
  \KOMAoptions{parskip=half}}
\makeatother
\usepackage{xcolor}
\IfFileExists{xurl.sty}{\usepackage{xurl}}{} % add URL line breaks if available
\IfFileExists{bookmark.sty}{\usepackage{bookmark}}{\usepackage{hyperref}}
\hypersetup{
  hidelinks,
  pdfcreator={LaTeX via pandoc}}
\urlstyle{same} % disable monospaced font for URLs
\usepackage{longtable,booktabs}
% Correct order of tables after \paragraph or \subparagraph
\usepackage{etoolbox}
\makeatletter
\patchcmd\longtable{\par}{\if@noskipsec\mbox{}\fi\par}{}{}
\makeatother
% Allow footnotes in longtable head/foot
\IfFileExists{footnotehyper.sty}{\usepackage{footnotehyper}}{\usepackage{footnote}}
\makesavenoteenv{longtable}
\setlength{\emergencystretch}{3em} % prevent overfull lines
\providecommand{\tightlist}{%
  \setlength{\itemsep}{0pt}\setlength{\parskip}{0pt}}
\setcounter{secnumdepth}{-\maxdimen} % remove section numbering

\author{}
\date{}

\begin{document}

\hypertarget{topological-constraint-on-the-electromagnetic-coupling-scale-from-a-discrete-ell1-obstruction}{%
\section{\texorpdfstring{Topological Constraint on the Electromagnetic
Coupling Scale from a Discrete \(\ell^1\)
Obstruction}{Topological Constraint on the Electromagnetic Coupling Scale from a Discrete \textbackslash ell\^{}1 Obstruction}}\label{topological-constraint-on-the-electromagnetic-coupling-scale-from-a-discrete-ell1-obstruction}}

\textbf{Author:} Jeremy H. Carroll \textbf{Version:} v1.0 (Pre-Print)
\textbf{Date:} March 2026

\begin{center}\rule{0.5\linewidth}{0.5pt}\end{center}

\hypertarget{abstract}{%
\subsection{Abstract}\label{abstract}}

We present a topological constraint on the electromagnetic coupling
scale derived from a single geometric seed: a sign contradiction on a
minimal graph. The \(\ell^1\) coboundary norm measures the
contradiction. Subdivision heals it into a 3-node lattice. From this
lattice, we derive the discrete Fourier transform as the unique
diagonalizing basis, the Born rule as Parseval's identity, and the
cyclic shift operator as the transport mechanism. The transport
amplitude across one edge of the resulting 2-simplex is shown to be
uniquely determined: the measure is Haar (by uniqueness of
shift-invariant measures), the observable is \(\cos\theta\) (by
linearity and symmetry constraints), and the domain is
\([-\pi/6, \pi/6]\) (by Voronoi bisection of the eigenvalue partition).
The resulting amplitude \(a = \sin(\pi/6)/(\pi/6) = 3/\pi\) yields a
unitarity deficit \(D = 1 - 9/\pi^2\), from which a candidate bare
coupling scale follows. The Casimir bridge to \(SU(3)\) and one-loop
renormalization group flow produce a candidate UV coupling consistent
with \(\alpha_{\mathrm{em}}^{-1}(0) \approx 137\) at the infrared scale.
A formal representation theorem, grounded in the Riesz--Markov--Kakutani
theorem, establishes that the operator class is not sampled but
exhausted. The derivation establishes uniqueness within the admissible
operator class defined by \(\ell^1\) stability; whether physical reality
is strictly restricted to this class remains an open question. The
derivation chain contains no free parameters and no imported physics
prior to the renormalization group stage. This work derives necessity,
not completeness: it shows what a minimal consistency condition forces,
not that it accounts for all of physics.

\begin{center}\rule{0.5\linewidth}{0.5pt}\end{center}

\hypertarget{introduction}{%
\subsection{1. Introduction}\label{introduction}}

The fine structure constant \(\alpha_{\mathrm{em}} \approx 1/137.036\)
has resisted first-principles derivation since its identification by
Sommerfeld in 1916. It enters the Standard Model as a free parameter,
and no existing framework--including string theory, loop quantum
gravity, or asymptotic safety--computes its value from internal
structure alone.

This work takes a different approach. Rather than beginning with a
quantum field theory and searching for a vacuum that reproduces
\(\alpha\), we begin with the minimal computable structure: a graph with
a sign contradiction. We then show that resolving this contradiction,
under the requirement that the resolution be stable, forces a unique
chain of algebraic and geometric structures that produces a constrained
candidate consistent with observed scaling. The only foundational
requirement is that local corrections to inconsistency must reduce, not
increase, global inconsistency. Everything else follows as theorem.

\textbf{Scope.} This work derives necessity, not completeness. We show
what a minimal consistency condition uniquely forces; we do not claim
that this single seed accounts for all of physics. The derivation
boundary--what is derived versus what is imported--is explicitly
delineated in Section 10.

The derivation proceeds in stages:

\begin{enumerate}
\def\labelenumi{\arabic{enumi}.}
\item
  \textbf{The seed} (Section 3): A 2-node graph with sections
  \(s(A) = +1\), \(s(B) = -1\) has \(\ell^1\) coboundary tension
  \(\|\delta s\|_1 = 2\). Subdivision heals this to a 3-node lattice.
\item
  \textbf{Wave mechanics} (Section 4): The cyclic shift operator on the
  3-node lattice is diagonalized uniquely by the discrete Fourier
  transform. Parseval's identity forces the Born rule \(P = |\psi|^2\).
\item
  \textbf{The transport operator} (Section 5): The eigenvalue partition
  of \(S^1\) defines Voronoi cells. Edge-sharing bisects each cell. The
  transport amplitude is the Haar average of \(\cos\theta\) over the
  edge sector \([-\pi/6, \pi/6]\), yielding \(a = 3/\pi\).
\item
  \textbf{Uniqueness} (Section 6): A representation theorem proves the
  operator space collapses to a single element under four requirements
  (bounded, linear, local, single-step), each derived from \(\ell^1\)
  stability axioms.
\item
  \textbf{The coupling constant} (Section 7): The unitarity deficit
  \(D = 1 - (3/\pi)^2\) and the Casimir bridge to \(SU(3)\) yield
  \(\alpha_{\mathrm{GUT}}^{-1} \approx 25.54\). Renormalization group
  flow to the infrared scale produces
  \(\alpha_{\mathrm{em}}^{-1} \approx 137\).
\item
  \textbf{Geometric Symmetry} (Section 8): The equilateral 2-simplex
  tiles \(\mathbb{R}^2\) uniquely as the exactly flat \(\{3,6\}\)
  lattice. This is presented as a suggestive geometric observation, not
  a cosmological model.
\end{enumerate}

Section 9 addresses the relation to holographic and multiverse
frameworks. Section 10 delineates what is derived from what is imported.
Section 11 discusses open problems.

\begin{center}\rule{0.5\linewidth}{0.5pt}\end{center}

\hypertarget{notation-and-conventions}{%
\subsection{2. Notation and
Conventions}\label{notation-and-conventions}}

\begin{longtable}[]{@{}ll@{}}
\toprule
\begin{minipage}[b]{0.40\columnwidth}\raggedright
Symbol\strut
\end{minipage} & \begin{minipage}[b]{0.54\columnwidth}\raggedright
Definition\strut
\end{minipage}\tabularnewline
\midrule
\endhead
\begin{minipage}[t]{0.40\columnwidth}\raggedright
\(\|\delta s\|_1\)\strut
\end{minipage} & \begin{minipage}[t]{0.54\columnwidth}\raggedright
\(\ell^1\) coboundary norm: \(\sum_{\text{edges}} |s(v) - s(w)|\)\strut
\end{minipage}\tabularnewline
\begin{minipage}[t]{0.40\columnwidth}\raggedright
\(N\)\strut
\end{minipage} & \begin{minipage}[t]{0.54\columnwidth}\raggedright
Number of lattice nodes (derived: \(N = 3\))\strut
\end{minipage}\tabularnewline
\begin{minipage}[t]{0.40\columnwidth}\raggedright
\(S\)\strut
\end{minipage} & \begin{minipage}[t]{0.54\columnwidth}\raggedright
Cyclic shift operator on \(N\) nodes\strut
\end{minipage}\tabularnewline
\begin{minipage}[t]{0.40\columnwidth}\raggedright
\(w_k\)\strut
\end{minipage} & \begin{minipage}[t]{0.54\columnwidth}\raggedright
Eigenvalue of \(S\): \(w_k = e^{2\pi i k/N}\)\strut
\end{minipage}\tabularnewline
\begin{minipage}[t]{0.40\columnwidth}\raggedright
\(\Omega\)\strut
\end{minipage} & \begin{minipage}[t]{0.54\columnwidth}\raggedright
Edge sector (Voronoi half-cell): \([-\pi/6, \pi/6]\)\strut
\end{minipage}\tabularnewline
\begin{minipage}[t]{0.40\columnwidth}\raggedright
\(a\)\strut
\end{minipage} & \begin{minipage}[t]{0.54\columnwidth}\raggedright
Transport amplitude: \(\mathrm{sinc}(\pi/6) = 3/\pi\)\strut
\end{minipage}\tabularnewline
\begin{minipage}[t]{0.40\columnwidth}\raggedright
\(D\)\strut
\end{minipage} & \begin{minipage}[t]{0.54\columnwidth}\raggedright
Unitarity deficit: \(1 - a^2 = 1 - 9/\pi^2\)\strut
\end{minipage}\tabularnewline
\begin{minipage}[t]{0.40\columnwidth}\raggedright
\(C_2(\mathrm{adj}), C_2(\mathrm{fund})\)\strut
\end{minipage} & \begin{minipage}[t]{0.54\columnwidth}\raggedright
Quadratic Casimir invariants of \(SU(3)\)\strut
\end{minipage}\tabularnewline
\bottomrule
\end{longtable}

\begin{center}\rule{0.5\linewidth}{0.5pt}\end{center}

\hypertarget{the-ell1-obstruction-and-its-resolution}{%
\subsection{\texorpdfstring{3. The \(\ell^1\) Obstruction and Its
Resolution}{3. The \textbackslash ell\^{}1 Obstruction and Its Resolution}}\label{the-ell1-obstruction-and-its-resolution}}

\hypertarget{the-minimal-contradiction}{%
\subsubsection{3.1 The Minimal
Contradiction}\label{the-minimal-contradiction}}

Consider the simplest non-trivial graph: two vertices \(A, B\) connected
by a single edge, with a section \(s: V \to \{+1, -1\}\) assigning
opposite signs. The \(\ell^1\) coboundary measures the disagreement:

\[\|\delta s\|_1 = |s(A) - s(B)| = 2\]

This is the minimal topological obstruction: \(H^1 \neq 0\) on the
graph. The section cannot be made globally consistent.

\hypertarget{subdivision-healing}{%
\subsubsection{3.2 Subdivision Healing}\label{subdivision-healing}}

The obstruction is resolved by iterative subdivision. When the tension
across an edge exceeds a threshold (set to 1), a new vertex is inserted
at the midpoint with interpolated value. Starting from \([+1, -1]\):

\begin{itemize}
\tightlist
\item
  Iteration 1: \([+1, 0.0, -1]\). Tensions: \(|1 - 0| = 1\),
  \(|0 - (-1)| = 1\). Healed.
\end{itemize}

The process terminates at \(N = 3\) nodes. This is deterministic: the
same seed always produces the same lattice. The subdivision threshold is
fixed by normalization of the \(\ell^1\) defect unit; any rescaling
\(\tau \mapsto \lambda\tau\) is equivalent to a change of units and does
not alter the fixed point \(N = 3\).

\hypertarget{why-ell1-and-not-another-norm}{%
\subsubsection{\texorpdfstring{3.3 Why \(\ell^1\) and Not Another
Norm}{3.3 Why \textbackslash ell\^{}1 and Not Another Norm}}\label{why-ell1-and-not-another-norm}}

The foundational requirement is not the \(\ell^1\) norm itself, but the
existence of a stable, monotone defect diagnostic---one under which
every local repair monotonically reduces the global obstruction measure.
The \(\ell^1\) coboundary norm is the unique structure satisfying this
requirement.

\textbf{Formal result.} Any norm on \(C^1(G, \mathbb{R})\) allowing
cancellation between positive and negative edge defects cannot guarantee
monotone defect detection on frustrated cycles. Specifically: let
\(\delta_e = s(v) - s(w)\) be the edge defect. A norm \(\|\cdot\|\) that
satisfies \(\|x + y\| < \|x\| + \|y\|\) for some sign-opposing \(x, y\)
(strict triangle inequality, as in \(\ell^2\)) necessarily admits
configurations where a single edge flip reduces the norm while
increasing the number of frustrated faces. Under \(\ell^1\), every edge
flip that reduces a local defect reduces the global diagnostic
monotonically, because \(\|\cdot\|_1 = \sum |\delta_e|\) decomposes as a
sum of independent non-negative terms.

This is established formally at three levels of generality in the
companion papers:

\begin{itemize}
\tightlist
\item
  \textbf{Paper 000} (Theorem 4.2, \(\ell^1\) Additive Universality):
  The \(\ell^1\) coboundary norm is the initial object in the category
  of faithful seminorms satisfying locality, additivity, symmetry, and
  non-degeneracy. Any such diagnostic is proportional to \(\ell^1\).
\item
  \textbf{Paper 001} (Theorem 2.2, Free Seminorm Classification): The
  result extends to Banach presheaf coboundaries under functoriality.
  The \(\ell^1\) norm is unique up to positive scalar among
  edge-additive, faithful, functorial seminorms.
\item
  \textbf{Paper 002} (Theorem 2.1, Cochain Classification):
  Coordinate-wise additivity over disjoint supports on finite graph
  cochains forces \(\ell^1\) geometry. This is the minimal combinatorial
  case.
\end{itemize}

Computationally, the stability argument is verified on \(N = 15\)
frustrated rings (\texttt{execute\_25}): \(\ell^1\) observers remain
stable while \(\ell^2\) and \(\ell^\infty\) observers drift past
stability thresholds. The norm inequality
\(\|\cdot\|_1 \geq \|\cdot\|_2 \geq \|\cdot\|_\infty\) ensures
\(\ell^1\) is maximally sensitive.

\begin{center}\rule{0.5\linewidth}{0.5pt}\end{center}

\hypertarget{wave-mechanics-from-graph-structure}{%
\subsection{4. Wave Mechanics from Graph
Structure}\label{wave-mechanics-from-graph-structure}}

\hypertarget{the-discrete-fourier-transform-as-unique-diagonalization}{%
\subsubsection{4.1 The Discrete Fourier Transform as Unique
Diagonalization}\label{the-discrete-fourier-transform-as-unique-diagonalization}}

The \(N\)-node lattice from Section 3 admits a cyclic shift operator
\(S\) (the translation-by-one-site map). The constraint that independent
subsystems factorize under this shift requires diagonalization. The
unique basis diagonalizing the cyclic convolution on \(N\) nodes is the
discrete Fourier transform (\texttt{execute\_09}). This is a theorem of
representation theory: the irreducible representations of the cyclic
group \(\mathbb{Z}_N\) are exactly the characters
\(e^{2\pi i k n / N}\).

Alternative bases (real-valued wavelets, random orthogonal matrices,
identity basis) are tested and fail to recover the convolution structure
(\texttt{execute\_09\_02}).

\hypertarget{the-born-rule-as-parsevals-identity}{%
\subsubsection{4.2 The Born Rule as Parseval's
Identity}\label{the-born-rule-as-parsevals-identity}}

The DFT satisfies Parseval's identity:
\(\sum_n |s_n|^2 = \sum_k |\hat{s}_k|^2\). An observer who lacks access
to the global \(\ell^1\) phase origin (gauge ignorance) must integrate
over this unknown phase. The unique translation-invariant measure on
\(U(1)\) is the Haar measure, and the resulting expectation value is
\(|\psi|^2\).

The Born rule is therefore not an axiom of the framework. It is a
theorem, derived at \texttt{execute\_11} from the DFT (which is derived
at \texttt{execute\_09} from the lattice, which is derived at
\texttt{execute\_00} from the obstruction). Probability enters the
derivation chain downstream of the seed, not before it.

\hypertarget{complex-phase-necessity}{%
\subsubsection{4.3 Complex Phase
Necessity}\label{complex-phase-necessity}}

The shift operator \(S\) on a directed cycle is asymmetric
(non-symmetric real matrix). By the spectral theorem, its eigenvalues
are generically complex. The complex plane is not imported; it is forced
by the directedness of causal transport on the lattice
(\texttt{execute\_10}). Real-valued diagonalization fails for asymmetric
circulant matrices.

\begin{center}\rule{0.5\linewidth}{0.5pt}\end{center}

\hypertarget{the-transport-operator-and-the-sinc-kernel}{%
\subsection{5. The Transport Operator and the Sinc
Kernel}\label{the-transport-operator-and-the-sinc-kernel}}

\hypertarget{eigenvalue-partition-of-s1}{%
\subsubsection{\texorpdfstring{5.1 Eigenvalue Partition of
\(S^1\)}{5.1 Eigenvalue Partition of S\^{}1}}\label{eigenvalue-partition-of-s1}}

The eigenvalues \(w_k = e^{2\pi i k / N}\) for \(k = 0, 1, \ldots, N-1\)
partition the unit circle into \(N\) equal arcs. For \(N = 3\), the
eigenangles are \(0, 2\pi/3, 4\pi/3\).

The Voronoi cell of each eigenvalue is the arc of angles closer to it
than to any other. For \(N\) equally-spaced points, each cell has width
\(2\pi/N\). The cell of \(w_0 = 1\) is \(\theta \in [-\pi/3, \pi/3]\).

\hypertarget{edge-sharing-bisection-and-the-locality-constraint}{%
\subsubsection{5.2 Edge-Sharing Bisection and the Locality
Constraint}\label{edge-sharing-bisection-and-the-locality-constraint}}

The eigenvalue partition defines the full Voronoi cell
\([-\pi/3, \pi/3]\). However, an edge is not a vertex. Transport is
fundamentally an edge-local operation. In a cycle graph \(C_3\), each
vertex is incident to exactly two edges (degree \(d = 2\)).

If we assigned the full vertex Voronoi cell \([-\pi/3, \pi/3]\) to both
incident edges, the edges would double-count the phase space surrounding
the vertex. This explicitly violates the strict independence and
additivity of the \(\ell^1\) norm. To preserve \(\ell^1\) additivity and
avoid double counting, the vertex measure must be partitioned equally
across its incident edges.

For degree 2, this forces the unique edge-local domain:
\[h = \frac{\pi/3}{1} = \frac{\pi}{3}, \quad \text{half-sector} = \frac{h}{2} = \frac{\pi}{6}\]

The edge sector \(\Omega = [-\pi/6, \pi/6]\) is therefore not a
geometric coincidence. Any domain larger than \([-\pi/6, \pi/6]\)
double-counts vertex measure; any smaller domain breaks symmetry. Thus
\([-\pi/6, \pi/6]\) is the unique admissible edge-local domain. This
partition is the unique measure-preserving decomposition of the vertex
Voronoi cell compatible with \(\ell^1\)-additive edge independence.

\hypertarget{the-forced-observable}{%
\subsubsection{5.3 The Forced Observable}\label{the-forced-observable}}

The transport operator \(T = e^{i\theta}\) is a unitary phase rotation.
To extract a real-valued amplitude compatible with the \(\ell^1\)
lattice (which evaluates real integer defects), we require the
observable \(f(\theta)\) to satisfy:

\begin{enumerate}
\def\labelenumi{\arabic{enumi}.}
\tightlist
\item
  \textbf{Linear in \(T\)}: \(f(T) = \mathrm{Re}(cT)\) for some \(c\),
  since \(T\) is the single-edge operator.
\item
  \textbf{Real-valued}: The \(\ell^1\) lattice is real.
\item
  \textbf{Normalized}: \(f(0) = 1\) (identity transport has unit
  amplitude).
\item
  \textbf{Even}: \(f(\theta) = f(-\theta)\) (the edge sector is
  symmetric about \(\theta = 0\)).
\end{enumerate}

Starting from \(f(\theta) = c \cdot e^{i\theta} + d \cdot e^{-i\theta}\)
with 4 real parameters, these constraints reduce: - Real \(\Rightarrow\)
2 parameters - Normalized \(\Rightarrow\) 1 parameter - Even
\(\Rightarrow\) 0 parameters

The unique survivor is
\(f(\theta) = \cos\theta = \mathrm{Re}(e^{i\theta})\).

\hypertarget{the-haar-measure}{%
\subsubsection{5.4 The Haar Measure}\label{the-haar-measure}}

The integration measure over \(\Omega\) must be shift-invariant (the
transport does not privilege any phase within the sector). By Haar's
theorem (1933), the unique such measure on a compact group is the
uniform (Lebesgue) measure.

\hypertarget{the-transport-amplitude}{%
\subsubsection{5.5 The Transport
Amplitude}\label{the-transport-amplitude}}

Combining the unique measure, observable, and domain:

\[a = \frac{1}{|\Omega|} \int_{-\pi/6}^{\pi/6} \cos\theta \, d\theta = \frac{\sin(\pi/6)}{\pi/6} = \frac{1/2}{\pi/6} = \frac{3}{\pi} \approx 0.9549\]

This is the sinc function evaluated at \(\pi/6\). The result contains no
free parameters.

\begin{center}\rule{0.5\linewidth}{0.5pt}\end{center}

\hypertarget{the-uniqueness-theorem}{%
\subsection{6. The Uniqueness Theorem}\label{the-uniqueness-theorem}}

\hypertarget{exhaustive-falsification}{%
\subsubsection{6.1 Exhaustive
Falsification}\label{exhaustive-falsification}}

Eight combinations of alternative observables (\(\cos\theta\),
\(\cos 2\theta\)) and domains (\([-\pi/12, \pi/12]\),
\([-\pi/6, \pi/6]\), \([-\pi/3, \pi/3]\), \([-\pi, \pi]\)) are tested.
Only \(\cos\theta\) over \([-\pi/6, \pi/6]\) yields \(3/\pi\). The
coincidence that \(\cos 2\theta\) over \([-\pi/12, \pi/12]\) also gives
\(3/\pi\) numerically is explained by the identity
\(\mathrm{sinc}(x) = \mathrm{sinc}(2 \cdot x/2)\), but this combination
is doubly excluded: \(\cos 2\theta = \mathrm{Re}(T^2)\) is quadratic in
\(T\) (violating linearity), and \([-\pi/12, \pi/12]\) has no
topological derivation.

\hypertarget{higher-harmonics-exclusion}{%
\subsubsection{6.2 Higher Harmonics
Exclusion}\label{higher-harmonics-exclusion}}

The observable \(\cos(n\theta) = \mathrm{Re}(T^n)\) corresponds to the
\(n\)-edge transport operator, not the single-edge operator. The
per-edge amplitude uses \(\langle k | S | k \rangle = w_k\), not
\(\langle k | S^n | k \rangle = w_k^n\). Higher harmonics conflate
multi-edge holonomy with single-edge transport.

Computationally: the average
\(\langle \cos(n\theta) \rangle_\Omega \neq a^n\) for \(n > 1\),
confirming these are distinct physical quantities.

\hypertarget{the-representation-theorem}{%
\subsubsection{6.3 The Representation
Theorem}\label{the-representation-theorem}}

\textbf{Theorem.} Let \(T: \mathcal{H} \to \mathbb{R}\) be a transport
functional satisfying: - \textbf{(R1) Bounded}:
\(|T[\psi]| \leq C \|\psi\|\) for some finite \(C\). - \textbf{(R2)
Linear}:
\(T[\alpha\psi_1 + \beta\psi_2] = \alpha T[\psi_1] + \beta T[\psi_2]\).
- \textbf{(R3) Local}: \(T\) depends only on the state within the edge
sector \(\Omega\). - \textbf{(R4) Single-step}: \(T\) measures the
amplitude for one edge traversal.

Then \(T\) is uniquely representable as:

\[T[\psi] = \int_\Omega \psi(\theta) \, d\mu(\theta)\]

with \$\mu = \$ Haar measure, \(\psi = \cos\theta\),
\(\Omega = [-\pi/6, \pi/6]\).

\textbf{Proof sketch.} 1. (R1) + (R2) imply \(T\) is a bounded linear
functional on \(C(\Omega)\), the space of continuous functions over the
compact domain \(\Omega\). By the Riesz--Markov--Kakutani representation
theorem (1909/1938/1941), \(T\) is representable as an integral against
a unique regular Borel measure. 2. This integral has the form
\((\text{measure}, \text{integrand}, \text{domain}) = (M, O, D)\). 3.
\(M\) is unique by Haar's theorem (Section 5.4). 4. \(O\) is unique by
the constraint analysis of Section 5.3. 5. \(D\) is unique by the
Voronoi bisection of Section 5.2.

\textbf{Escape routes closed:} - \emph{Nonlocal kernels}: violate (R3).
The Voronoi sectors are disjoint; cross-sector coupling means transport
across edge \(e\) depends on states at distant edge \(e'\). -
\emph{Distributional operators}: the transport amplitude is a scalar in
\([0, 1]\), not a distribution. - \emph{Nonlinear functionals}: violate
(R2), which is forced by the DFT-derived Hilbert space structure. -
\emph{Path-dependent operators}: violate (R4). The shift operator \(S\)
is Markov; multi-step transport decomposes as iterated single-step. -
\emph{Full-support weighted operators}: Haar uniqueness forces uniform
measure; only \(h = \pi/6\) yields \(3/\pi\).

\hypertarget{axiom-justification}{%
\subsubsection{6.4 Axiom Justification}\label{axiom-justification}}

The requirements R1--R4 are not external assumptions. They are the
operator-language translations of four \(\ell^1\) stability axioms, each
proven independently necessary (\texttt{execute\_24}):

\begin{longtable}[]{@{}lll@{}}
\toprule
\begin{minipage}[b]{0.32\columnwidth}\raggedright
Operator Requirement\strut
\end{minipage} & \begin{minipage}[b]{0.23\columnwidth}\raggedright
\(\ell^1\) Axiom\strut
\end{minipage} & \begin{minipage}[b]{0.37\columnwidth}\raggedright
Failure Mode if Violated\strut
\end{minipage}\tabularnewline
\midrule
\endhead
\begin{minipage}[t]{0.32\columnwidth}\raggedright
R1 Bounded\strut
\end{minipage} & \begin{minipage}[t]{0.23\columnwidth}\raggedright
Faithfulness\strut
\end{minipage} & \begin{minipage}[t]{0.37\columnwidth}\raggedright
Unbounded \(\Rightarrow\) unfaithful diagnostic\strut
\end{minipage}\tabularnewline
\begin{minipage}[t]{0.32\columnwidth}\raggedright
R2 Linear\strut
\end{minipage} & \begin{minipage}[t]{0.23\columnwidth}\raggedright
Coproduct compatibility\strut
\end{minipage} & \begin{minipage}[t]{0.37\columnwidth}\raggedright
Nonlinear \(\Rightarrow\) defect cancellation\strut
\end{minipage}\tabularnewline
\begin{minipage}[t]{0.32\columnwidth}\raggedright
R3 Local\strut
\end{minipage} & \begin{minipage}[t]{0.23\columnwidth}\raggedright
Locality\strut
\end{minipage} & \begin{minipage}[t]{0.37\columnwidth}\raggedright
Nonlocal \(\Rightarrow\) missed distributed defects\strut
\end{minipage}\tabularnewline
\begin{minipage}[t]{0.32\columnwidth}\raggedright
R4 Single-step\strut
\end{minipage} & \begin{minipage}[t]{0.23\columnwidth}\raggedright
Contractive monotonicity\strut
\end{minipage} & \begin{minipage}[t]{0.37\columnwidth}\raggedright
Multi-step \(\Rightarrow\) amplifying corrections\strut
\end{minipage}\tabularnewline
\bottomrule
\end{longtable}

The \(\ell^1\) norm itself is forced by observer stability (Section
3.3). Therefore: the derivation chain from obstruction to \(\alpha\)
contains no free axioms. Every requirement is derived from the one
before it.

\begin{center}\rule{0.5\linewidth}{0.5pt}\end{center}

\hypertarget{the-coupling-constant}{%
\subsection{7. The Coupling Constant}\label{the-coupling-constant}}

\hypertarget{the-unitarity-deficit}{%
\subsubsection{7.1 The Unitarity Deficit}\label{the-unitarity-deficit}}

The transport amplitude \(a = 3/\pi < 1\) implies that traversing one
edge of the 2-simplex is contractive. The unitarity deficit per edge is:

\[D = 1 - a^2 = 1 - \frac{9}{\pi^2} \approx 0.08860\]

For the full triangular loop, the operator \(T_\triangle = T_1 T_2 T_3\)
has spectrum
\(\mathrm{Spec}(T_\triangle^\dagger T_\triangle) = \{a^6, a^6, a^6\}\),
and the scalar deficit \(D_\triangle = 1 - a^6 \approx 0.2417\) is
gauge-invariant under \(SU(3)\) transformations (\texttt{execute\_33}).

\hypertarget{the-casimir-bridge}{%
\subsubsection{7.2 The Casimir Bridge}\label{the-casimir-bridge}}

The 2-simplex has 3 oriented edges, making the boundary cochain
\(C^1 \cong \mathbb{C}^3\). The cyclic transport generates a
\(\mathbb{Z}_3\) center, and \(\det = 1\) (conservation). By Lie algebra
classification, these constraints strongly single out \(SU(3)\) as the
minimal consistent candidate.

\hypertarget{uniqueness-of-su3}{%
\paragraph{\texorpdfstring{7.2.1 Uniqueness of
\(SU(3)\)}{7.2.1 Uniqueness of SU(3)}}\label{uniqueness-of-su3}}

A referee may ask why not \(SU(2) \times U(1)\) or a larger group. The
transport structure acts on a 3-dimensional complex cochain space
(\(C^1 \cong \mathbb{C}^3\)) with a \(\mathbb{Z}_3\) phase symmetry and
a global conservation law. The representation must therefore satisfy all
of the following geometrically-derived constraints:

\begin{enumerate}
\def\labelenumi{\arabic{enumi}.}
\tightlist
\item
  \textbf{Faithful representation on \(\mathbb{C}^3\)}: The group must
  act faithfully on the 3 oriented complex edges. This excludes groups
  whose lowest faithful representation is higher-dimensional, and
  excludes \(SU(2)\) which lacks a 3D fundamental representation. It
  also excludes \(SO(3)\) which is real, not complex.
\item
  \textbf{Center containing \(\mathbb{Z}_3\)}: The cyclic shift \(S\) on
  3 nodes enforces eigenvalues containing \(e^{2\pi i/3}\). The group
  center must contain this symmetry. This excludes \(SU(2)\) (center
  \(\mathbb{Z}_2\)) and \(Sp(n)\) (center \(\mathbb{Z}_2\)).
\item
  \textbf{Determinant-preserving (\(\det = 1\))}: The \(\ell^1\)
  subdivision healing explicitly conserves the total section mass,
  translating to a determinant of 1. This excludes \(U(3)\) (which has
  \(\det \in U(1)\)) and \(GL(3)\) (not unitary).
\item
  \textbf{Non-abelian}: The 2-simplex enforces correlated transport
  across edges to close the loop holonomy (\(S^3 = I\)). An abelian
  product like \(U(1)^3\) cannot enforce edge-to-edge correlation and
  curvature.
\item
  \textbf{Irreducible}: The 2-simplex is a single, indecomposable
  connected face. Realizing the group via a block-diagonal (reducible)
  representation like \(SU(2) \times U(1)\) violates this geometric
  indecomposability.
\end{enumerate}

By Cartan's classification of compact Lie groups, these constraints
leave exactly one minimal admissible solution: \(SU(3)\). Larger groups
would introduce extraneous generators and degrees of freedom unsupported
by the exact dimension of the 2-simplex topology. \(SU(3)\) is not just
mathematically natural; it is uniquely forced as the minimal closed
solution.

The mapping from edge transport to face curvature is not an arbitrary
choice---it corresponds to the the coboundary lift \(d^1: C^1 \to C^2\).
The fundamental defining action resides on the edges (the fundamental
representation, \(\dim 3\)), while the closed loop holonomy measures
curvature, taking values in the Lie algebra (the adjoint representation,
\(\dim 8\)).

When translating the amplitude deficit \(D\) from the fundamental
representation on the edge to the adjoint representation on the face,
the geometric scaling factor is exactly the ratio of their quadratic
Casimir invariants. This scaling is a theorem of Lie algebra
representations dictating how traces scale when lifting an algebraic
structure from its fundamental defining space to its adjoint bracket
space. It is not an imported QFT axiom, but a structural scaling ratio:

\[\frac{1}{\alpha_{\mathrm{GUT}}} = \frac{C_2(\mathrm{adj})}{C_2(\mathrm{fund})} \times \frac{1}{D} = \frac{3}{4/3} \times \frac{\pi^2}{\pi^2 - 9} = \frac{9}{4} \times \frac{1}{D} \approx 25.54\]

\hypertarget{renormalization-group-flow}{%
\subsubsection{7.3 Renormalization Group
Flow}\label{renormalization-group-flow}}

\textbf{Note on imported quantities.} The RG flow from
\(\alpha_{\mathrm{GUT}}^{-1} \approx 25.54\) to the infrared scale
requires beta-function coefficients, which encode the matter content of
the theory. For the MSSM (\(N_g = 3\), \(N_h = 2\)), the one-loop
coefficients are \(\beta_1 = 33/5\), \(\beta_2 = 1\), \(\beta_3 = -3\).

These coefficients are not derived from the \(\ell^1\) obstruction in
this work. They represent the interface between the topological
derivation (Sections 3--6) and phenomenology. The derivation of gauge
group representations from \(\ell^1\) topology is allocated to the Grand
Unification program (Paper 009).

The derivation provides a candidate UV coupling
\(\alpha_{\mathrm{GUT}}^{-1} \approx 25.54\) consistent with
\(\alpha_{\mathrm{em}}\) under RG flow; the IR value
\(\alpha_{\mathrm{em}}^{-1} \approx 137\) depends on matter content,
which remains to be derived from \(\ell^1\) topology in subsequent work.

With this caveat, the one-loop RG flow yields:

\[\alpha_{\mathrm{em}}^{-1}(M_Z) \approx 128, \qquad \alpha_{\mathrm{em}}^{-1}(0) \approx 137\]

consistent with experimental measurement.

\begin{center}\rule{0.5\linewidth}{0.5pt}\end{center}

\hypertarget{the-infinite-tiling-and-geometric-symmetry}{%
\subsection{8. The Infinite Tiling and Geometric
Symmetry}\label{the-infinite-tiling-and-geometric-symmetry}}

\hypertarget{the-36-tiling}{%
\subsubsection{\texorpdfstring{8.1 The \(\{3,6\}\)
Tiling}{8.1 The \textbackslash\{3,6\textbackslash\} Tiling}}\label{the-36-tiling}}

The 2-simplex from Section 3 is equilateral (forced by \(\mathbb{Z}_3\)
edge equivalence, \texttt{execute\_29}). Equilateral triangles tile the
Euclidean plane uniquely as the \(\{3,6\}\) Schlafli tiling: 6 triangles
meet at each vertex.

The Regge deficit at every vertex is:

\[\epsilon = 2\pi - 6 \times \frac{\pi}{3} = 0\]

This is zero identically, not approximately. The tiling is exactly flat.

\hypertarget{a-suggestive-geometric-observation}{%
\subsubsection{8.2 A Suggestive Geometric
Observation}\label{a-suggestive-geometric-observation}}

The \(\{3,6\}\) tiling identifies an exactly flat underlying continuous
structure. While an initially perfectly flat discrete geometry is
structurally encouraging given observed cosmic flatness, we explicitly
classify this as a suggestive geometric observation, not a cosmological
model. It does not contain dynamics, expansion, or perturbation
modeling, and determining whether this fundamentally relates to the
physical universe's flatness problem lies strictly outside the scope of
this paper's claims.

\begin{center}\rule{0.5\linewidth}{0.5pt}\end{center}

\hypertarget{relation-to-alternative-frameworks}{%
\subsection{9. Relation to Alternative
Frameworks}\label{relation-to-alternative-frameworks}}

\hypertarget{prior-probability-and-the-landscape}{%
\subsubsection{9.1 Prior Probability and the
Landscape}\label{prior-probability-and-the-landscape}}

Frameworks that begin with a probability distribution over possible
universes (the string landscape, eternal inflation with a measure
problem, Tegmark's Level IV multiverse) presuppose a mathematical
structure---probability---that is not available at the foundational
level of this derivation.

In the \(\ell^1\) pipeline, the Born rule \(P = |\psi|^2\) is derived at
Section 4.2 from Parseval's identity, which requires the DFT (Section
4.1), which requires the \(N\)-node lattice (Section 3.2), which
requires the obstruction (Section 3.1). Probability enters the
derivation chain downstream of the seed. Any framework that invokes
probability at the foundational level uses a tool that does not yet
exist at the point where the selection would occur.

\hypertarget{the-everett-interpretation}{%
\subsubsection{9.2 The Everett
Interpretation}\label{the-everett-interpretation}}

Everett's many-worlds interpretation (MWI) requires three elements: (E1)
a Hilbert space, (E2) unitary evolution, and (E3) branch weights given
by the \(\ell^2\) norm.

The \(\ell^1\) pipeline derives (E1) at Section 4.1 and (E2) as an
emergent infrared limit. However, it denies (E3): the \(\ell^2\) norm is
not fundamental but derived from \(\ell^1\) via the gauge-ignorance
functor (Parseval's identity). Observers using \(\ell^2\) defect
accounting are unstable on frustrated graphs (\texttt{execute\_25}), and
the \(\ell^1\) healing is deterministic with a unique fixed point
(\(N = 3\)).

MWI correctly identifies the Hilbert space but incorrectly elevates the
derived \(\ell^2\) norm to foundational status. Everett's branches are
epistemic (observer-relative projections of a single \(\ell^1\)
structure), not ontic (independently existing universes).

\hypertarget{holography-and-adscft}{%
\subsubsection{9.3 Holography and AdS/CFT}\label{holography-and-adscft}}

The AdS/CFT correspondence and related holographic frameworks operate
within a continuous \(L^2\) Hilbert space and explore a broad class of
consistent quantum field theories and gravitational duals. These
approaches have achieved remarkable success in characterizing strongly
coupled systems and establishing deep mathematical dualities.

The present work takes a different foundational approach. Rather than
assuming a Hilbert structure, it derives the underlying space from a
discrete \(\ell^1\) topological obstruction. The focus is not on
classifying all consistent theories, but on determining whether a
minimal consistency condition uniquely fixes physical structure.

Several structural features distinguish the two approaches:

\begin{enumerate}
\def\labelenumi{\arabic{enumi}.}
\item
  \textbf{Geometry}: The \(\ell^1\) tiling yields exactly flat
  (\(\epsilon = 0\)) Euclidean geometry. AdS/CFT requires negative
  cosmological constant (negative curvature). The pipeline forces vertex
  valence \(k = 2\pi/(\pi/3) = 6\), yielding the \(\{3,6\}\) tiling. The
  hyperbolic \(\{3,7\}\) tiling required for AdS-like geometry is not
  accessible.
\item
  \textbf{Discreteness}: The holographic principle (Bekenstein bound,
  Ryu--Takayanagi formula) presupposes a smooth manifold with
  well-defined areas. The \(\ell^1\) lattice is fundamentally discrete;
  information content is the coboundary norm on edges, not an
  area-entropy relationship.
\item
  \textbf{Information structure}: In holographic frameworks, information
  on a \((d-1)\)-dimensional boundary encodes a \(d\)-dimensional bulk.
  In the \(\ell^1\) framework, information lives on edges (1-cells).
  There is no higher-dimensional bulk for a lower-dimensional boundary
  to encode.
\item
  \textbf{Methodology}: AdS/CFT is a duality relating two descriptions
  of the same physics, both assuming extensive prior structure. The
  \(\ell^1\) pipeline is a derivation from a single seed. These are
  complementary approaches: holographic frameworks characterize what is
  \emph{possible}; this construction investigates what is
  \emph{necessary}.
\end{enumerate}

\begin{center}\rule{0.5\linewidth}{0.5pt}\end{center}

\hypertarget{derivation-boundary-what-is-derived-and-what-is-imported}{%
\subsection{10. Derivation Boundary: What Is Derived and What Is
Imported}\label{derivation-boundary-what-is-derived-and-what-is-imported}}

Intellectual honesty requires a clear delineation.

\hypertarget{derived-from-the-ell1-obstruction-no-external-input}{%
\subsubsection{\texorpdfstring{10.1 Derived from the \(\ell^1\)
Obstruction (No External
Input)}{10.1 Derived from the \textbackslash ell\^{}1 Obstruction (No External Input)}}\label{derived-from-the-ell1-obstruction-no-external-input}}

\begin{longtable}[]{@{}lll@{}}
\toprule
Structure & Source & Section\tabularnewline
\midrule
\endhead
\(N = 3\) lattice & Subdivision healing & 3.2\tabularnewline
DFT as unique basis & Representation theory of \(\mathbb{Z}_N\) &
4.1\tabularnewline
Born rule \(P = |\psi|^2\) & Parseval's identity & 4.2\tabularnewline
Complex phase & Asymmetric shift operator & 4.3\tabularnewline
Voronoi partition & Metric geometry on \(S^1\) & 5.1\tabularnewline
Edge sector \([-\pi/6, \pi/6]\) & Bisection by vertex degree &
5.2\tabularnewline
Observable \(\cos\theta\) & Linearity + symmetry constraints &
5.3\tabularnewline
Haar measure & Uniqueness of shift-invariant measure &
5.4\tabularnewline
\(a = 3/\pi\) & Sinc kernel & 5.5\tabularnewline
\(D = 1 - 9/\pi^2\) & Unitarity deficit & 7.1\tabularnewline
\(SU(3)\) gauge group & Lie algebra classification & 7.2\tabularnewline
\(\alpha_{\mathrm{GUT}}^{-1} \approx 25.54\) & Casimir scaling &
7.2\tabularnewline
Euclidean flatness & Regge calculus on \(\{3,6\}\) & 8.1\tabularnewline
Representation theorem & Riesz--Markov--Kakutani + \(\ell^1\) axioms &
6.3--6.4\tabularnewline
\bottomrule
\end{longtable}

\hypertarget{imported-from-phenomenology}{%
\subsubsection{10.2 Imported from
Phenomenology}\label{imported-from-phenomenology}}

\begin{longtable}[]{@{}lll@{}}
\toprule
Structure & Status & Section\tabularnewline
\midrule
\endhead
MSSM beta-function coefficients & Input (matter content) &
7.3\tabularnewline
Couplings at \(M_Z\) & Experimental boundary condition &
7.3\tabularnewline
\(M_{\mathrm{Planck}}\) & External gravitational scale &
7.3\tabularnewline
\(U(1) \times SU(2)\) gauge structure & Not derived from \(\ell^1\) in
this work & 7.3\tabularnewline
\bottomrule
\end{longtable}

\hypertarget{honest-assessment}{%
\subsubsection{10.3 Honest Assessment}\label{honest-assessment}}

The genuine achievement of this work is the derivation of \(a = 3/\pi\)
and \(D = 1 - 9/\pi^2\) from a single topological seed with zero free
parameters (Sections 3--6). The Casimir bridge to \(SU(3)\) (Section
7.2) is structurally clean. The RG flow to \(1/137\) (Section 7.3)
requires phenomenological inputs that are not yet derived from
\(\ell^1\) topology.

The derivation of the full Standard Model gauge structure
(\(U(1) \times SU(2) \times SU(3)\), generation count, Higgs sector)
from the \(\ell^1\) obstruction is an open problem allocated to the
Grand Unification program (Paper 009).

\begin{center}\rule{0.5\linewidth}{0.5pt}\end{center}

\hypertarget{open-problems}{%
\subsection{11. Open Problems}\label{open-problems}}

\begin{enumerate}
\def\labelenumi{\arabic{enumi}.}
\item
  \textbf{Derive the electroweak sector.} The \(SU(3)\) gauge group is
  forced by the 2-simplex structure. The \(U(1) \times SU(2)\) sector
  and its breaking pattern are not yet derived from \(\ell^1\) topology.
\item
  \textbf{Derive the beta-function coefficients.} The matter content (3
  generations, 2 Higgs doublets) enters as an input to the RG flow. A
  first-principles derivation would eliminate the last imported
  quantities.
\item
  \textbf{Derive gravity.} The Regge deficit is identically zero on the
  \(\{3,6\}\) tiling. Curvature (gravity) requires defects in the
  tiling. The mechanism by which \(\ell^1\) defects source geometric
  curvature is not yet formalized.
\item
  \textbf{Experimental predictions.} The framework predicts that
  Planck-scale physics is non-unitary (Section 4.3), with unitarity
  emerging as an infrared limit. Observable consequences at accessible
  energies remain to be identified.
\item
  \textbf{The origin of distinction.} The derivation begins at the first
  computable condition: the existence of a sign contradiction on a
  graph. Why such a distinction exists lies outside the domain of
  physical theory and enters metaphysics. The pipeline is complete
  within the scope of computable structure.
\end{enumerate}

\begin{center}\rule{0.5\linewidth}{0.5pt}\end{center}

\hypertarget{conclusion}{%
\subsection{12. Conclusion}\label{conclusion}}

We have presented a computational derivation chain from a single
topological seed (a sign contradiction on a minimal graph) to a
topological constraint on the electromagnetic coupling scale, passing
through the emergence of wave mechanics, the Born rule, complex phase,
the sinc transport kernel, and the \(SU(3)\) Casimir bridge.

The central result is the transport amplitude
\(a = \mathrm{sinc}(\pi/6) = 3/\pi\), proven unique by a representation
theorem grounding the operator classification in the
Riesz--Markov--Kakutani theorem, with requirements derived from
\(\ell^1\) stability axioms rather than imposed externally.

The derivation contains zero free parameters from seed to
\(\alpha_{\mathrm{GUT}}^{-1} \approx 25.54\). The RG flow to
\(\alpha_{\mathrm{em}}^{-1} \approx 137\) requires phenomenological
inputs (beta-functions, scale matching) that represent the current
interface between topological derivation and Standard Model
phenomenology.

Under the requirement of \(\ell^1\)-stable observer consistency, the
admissible operator space collapses to a single solution. The derivation
establishes uniqueness within the admissible operator class defined by
\(\ell^1\) stability; whether physical reality is strictly restricted to
this class remains an open question. This leads to a topological
constraint, in contrast to frameworks that admit large theory spaces or
rely on anthropic selection. This distinction reflects a difference in
foundational perspective: holographic and landscape approaches
characterize what is possible, while the present construction
investigates what is necessary.

The complete derivation chain, including all computational proofs and
falsification tests, is available in the companion script
\texttt{master\_l1\_derivation.py}.

\begin{center}\rule{0.5\linewidth}{0.5pt}\end{center}

\hypertarget{references}{%
\subsection{References}\label{references}}

\begin{enumerate}
\def\labelenumi{\arabic{enumi}.}
\tightlist
\item
  Haar, A. (1933). Der Massbegriff in der Theorie der kontinuierlichen
  Gruppen. \emph{Annals of Mathematics}, 34(1), 147--169.
\item
  Riesz, F. (1909). Sur les operations fonctionnelles lineaires.
  \emph{Comptes Rendus}, 149, 974--977.
\item
  Kakutani, S. (1941). Concrete representation of abstract (M)-spaces.
  \emph{Annals of Mathematics}, 42(4), 994--1024.
\item
  Regge, T. (1961). General relativity without coordinates. \emph{Il
  Nuovo Cimento}, 19(3), 558--571.
\item
  Maldacena, J. (1999). The large-\(N\) limit of superconformal field
  theories and supergravity. \emph{International Journal of Theoretical
  Physics}, 38(4), 1113--1133.
\item
  Everett, H. (1957). ``Relative state'' formulation of quantum
  mechanics. \emph{Reviews of Modern Physics}, 29(3), 454--462.
\item
  Shannon, C. E. (1948). A mathematical theory of communication.
  \emph{Bell System Technical Journal}, 27(3), 379--423.
\item
  Carroll, J. H. (2026). \emph{Projection Obstruction Theory: Retraction
  Nonlinearity, \(\ell^1\) Rigidity, and Density Scaling.} (Paper 000).
  Zenodo. https://doi.org/10.5281/zenodo.19151803
\item
  Carroll, J. H. (2026). \emph{The Free \(\ell^1\) Seminorm on Banach
  Presheaf Coboundaries.} (Paper 001). Zenodo.
  https://doi.org/10.5281/zenodo.19152115
\item
  Carroll, J. H. (2026). \emph{Coordinate-Wise Additivity and the
  \(\ell^1\) Norm on Finite Graph Cochains.} (Paper 002). Zenodo.
  https://doi.org/10.5281/zenodo.19152599
\item
  Carroll, J. H. (2026). \emph{Hodge Structure and Gauge Equivalence of
  \(\ell^1\) Defect Fields.} (Paper 003). Zenodo.
  https://doi.org/10.5281/zenodo.19152799
\item
  Carroll, J. H. (2026). \emph{Universal Obstruction Theory: The
  \(\ell^1\) Topological Framework.} (Paper 004). Zenodo.
  https://doi.org/10.5281/zenodo.19154775
\item
  Carroll, J. H. (2026). \emph{Autopoietic Cohomology: Iterative
  Obstruction Repair on Causal Posets.} (Paper 005). Zenodo.
  https://doi.org/10.5281/zenodo.19155011
\end{enumerate}

\begin{center}\rule{0.5\linewidth}{0.5pt}\end{center}

\end{document}
